\chapter{引言}\label{ux5f15ux8a00}

\section{研究背景}\label{ux7814ux7a76ux80ccux666f}

建筑业作为关系国计民生的基础性产业，是国民经济的重要组成部分。据中国建筑业协会发布的《2025年建筑业发展统计分析》显示，2025全年全社会建筑业实现增加值86425.1亿元，占国内生产总值的6.2\%\cite{E78S579E}。在建筑施工作业中，塔式起重机是使用最广泛的垂直运输设备，具有起升高度大、作业范围广、工作效率高等优势\cite{G8HEXPSB,3LHPNGFB,CSVHSWVZ}，其吊装作业承担着钢筋笼、模板、物料等关键构件的吊运任务。然而，塔吊施工环境普遍存在空间狭窄、视野受限、动态干扰复杂等特点，加之运行高度大、工作半径长，一旦发生碰撞或人员误入吊装路径，后果往往极为严重。据统计，2014年至2022年间全国塔式起重机安全事故超过220起\cite{KWVSPS3W}，事故类型分布如表\ref{tab:accident_distribution}所示。其中，吊物伤人与物体碰撞（碰撞挤压）类事故累计超过75起，占比达34\%以上，是威胁吊装安全的主要事故类型。因此，实现吊装过程的实时环境感知与风险预警，是当前智慧工地安全管理的重要任务。

\begin{table}[htbp]
\centering
\caption{2014--2022年塔式起重机事故类型统计\cite{KWVSPS3W}}
\label{tab:accident_distribution}
\begin{tabular*}{\textwidth}{@{\extracolsep{\fill}}>{\centering\arraybackslash}m{2.5cm}>{\centering\arraybackslash}m{1.8cm}>{\centering\arraybackslash}m{2.0cm}>{\centering\arraybackslash}m{6.5cm}}
\Xhline{1pt}
\textbf{事故类型} & \textbf{起数} & \textbf{占比排序} & \textbf{典型致因} \\
\hline
\multicolumn{4}{l}{\textbf{吊装碰撞直接相关事故}} \\
\Xhline{0.5pt}
吊物伤人 & 54 & 第2位 & 吊装构件在吊装过程中撞击下方人员 \\
物体碰撞 & 21 & 第4位 & 吊物与周边结构、设备发生碰撞挤压 \\
\hline
\multicolumn{4}{l}{\textbf{其他事故类型}} \\
\Xhline{0.5pt}
倒塌 & 79 & 第1位 & 超载或摆动过大导致整机失稳 \\
高处坠落 & 52 & 第3位 & 作业人员从吊装平台或吊物上坠落 \\
其他 & 14 & 第5-7位 & 包括大臂折断、套架滑落、触电伤害等 \\
\Xhline{1pt}
\end{tabular*}
\end{table}

吊钩端是吊装过程中运动最频繁、最接近人员车辆及构件的部位，其近场作业具有以下显著特征：（1）空间运动高度不确定：吊钩随吊臂回转、钢丝绳伸缩及外力扰动持续发生三维空间运动，摆幅可达数米，传统固定式传感器难以覆盖其动态范围；（2）近场作业风险集中：施工人员在吊运构件下方或附近配合作业、车辆设备在近场区域穿行，均可能导致危险目标无意靠近吊钩端，叠加吊物摆动与构件碰撞等因素，易引发重大安全事故；（3）环境遮挡复杂且光照变化剧烈：钢筋笼、模板体系等大型构件会造成视野遮挡，而工地扬尘、夜间照明等不利因素对视觉系统影响显著；（4）传统方法难以实现实时近场避障：现有塔吊安全系统多依赖塔身安装的超声、雷达或视觉设备，检测范围难以覆盖吊钩周围高危近场区域，使关键危险点无法及时识别。

然而，现有塔吊安全感知研究多集中于塔身或塔臂端的宏观防碰撞与运行状态监测，例如多塔协同作业的塔臂干涉检测、塔机与周边建筑结构的距离监测等，而对实际吊装环节中风险最高的吊钩端区域直接感知关注显著不足\cite{9LK3VTEQ}。

需要指出的是，吊钩端避障的落脚点实质上是吊物避障。在吊装作业中，实际可能与周边结构或人员发生碰撞的对象是具有一定体积与复杂几何形态的吊装构件（如钢筋笼、模板、预制构件等），而非作为起吊连接节点的吊钩本体；吊物随钢丝绳摆动与塔臂回转呈现大幅空间位移与姿态变化，其动态外包络范围远超吊钩。因此，吊钩端感知的核心任务在于对吊物状态的准确估计，即通过感知手段获取吊物的三维位置、空间范围及运动趋势，而非仅对吊钩进行定位跟踪。然而，现有研究多将“吊钩端”简化为“吊钩位置”进行跟踪，缺乏对吊物本体空间状态的直接感知与估计方法。

从碰撞风险类型来看，吊钩端近场需防范的碰撞威胁可归纳为两类：一是静态空间侵入风险，即吊钩、吊物或吊具与周边静态结构（如建筑主体、脚手架、邻近塔吊、地面堆料等）之间的安全净空不足，尤其在吊物回转或摆动过程中，吊物轮廓与周边静态构件的距离持续变化，需要实时评估三维净空；二是动态交会碰撞风险，即施工人员、运输车辆、移动机械等动态危险目标进入吊钩端运动包络或吊装路径，与吊物形成短时交会。在基于塔身固定传感器的常规宏观防碰撞架构下，上述两类风险难以得到直接、精细的感知与量化评估。两类风险的预警逻辑存在本质差异：静态风险取决于当前及预测时刻吊物与周边结构的几何距离，动态风险则还需考虑目标相对运动趋势与碰撞时间裕度，这对传感器的近场覆盖范围、测距精度与实时处理能力提出了更高要求。

部分企业尝试在塔吊大臂端部署激光雷达与视觉传感器，以期通过远距离观测实现吊物周边的环境感知。然而，大臂端部署方案同样难以满足吊钩端的感知需求，受限于以下固有瓶颈：第一，在典型50\,m臂长工况下，大臂端传感器到吊物的距离通常超过40\,m，远距离观测导致LiDAR点云分辨率急剧下降，难以可靠检测人员、构件边缘等厘米级小目标，图像像素分辨率和测距精度亦难以满足近场安全净空的精确判定需求；第二，吊臂回转过程中传感器与吊物相对视角持续大幅变化，加之吊具与钢丝绳对远距离视线的频繁遮挡，导致大臂端传感器难以稳定跟踪吊物周边的动态危险目标；第三，大臂端传感器通常固定于塔吊结构，无法随吊钩端同步移动，其对吊钩正下方及侧方盲区覆盖存在难以克服的几何死角。上述局限表明，仅依靠塔臂端或大臂端感知方案无法有效覆盖吊钩端近场高危区域，将感知单元直接部署于吊钩端是从物理层面缩短传感距离、提升近场感知精度与持续性的必要手段。

这一部署思路的技术可行性，得益于近年来轻量化激光雷达和高性能嵌入式计算平台的快速发展。激光雷达具备高精度距离测量能力，能够在强光、弱光、粉尘等复杂工况中保持稳定输出，并可实现高达360°的水平视场（FoV），从而覆盖吊钩周围的全方位近场空间。相比之下，视觉传感器具有丰富的语义理解能力，特别是在目标识别与像素级分割方面具有显著优势，可为点云投影提供精确的语义边界约束。因此，构建“雷达为主、视觉为辅”的雷视融合感知架构能够充分结合两类传感器的互补特性，为实现塔吊吊钩端的实时、可靠避障预警系统创造了可行条件。

然而，在吊钩端直接部署感知方案同样面临独特的技术挑战。与固定于塔臂或塔身的传感器不同，吊钩端感知平台处于持续三维运动状态，传感器随吊臂回转、钢丝绳伸缩与外力扰动实时运动，观测基准随之动态改变。这一高动态特性对系统提出了两方面的核心难题：其一，一致性环境感知问题，在吊钩端自身持续运动条件下，传感器坐标系时刻变化，需要将不同时刻、不同视角的观测数据统一至一致的空间参考框架中，才能实现环境与吊物状态的准确估计；其二，动态预警策略设计问题，传统基于固定阈值的距离判据难以直接适用于运动基准下的净空评估，需要将传感器自身运动、目标相对运动与未来趋势纳入统一的风险建模框架。上述难题构成了吊钩端感知区别于固定式部署方案的核心技术难点，需要在感知架构与预警模型层面进行针对性设计。

基于上述背景，本研究面向吊钩端高风险作业场景，选用双激光雷达与工业相机设计并开发集成式吊钩端感知系统，以点云投影为技术手段实现吊物状态准确估计与近场三维环境的高精度感知，支撑吊装构件识别与碰撞风险判定，减少施工现场吊装作业中的碰撞风险、保障人员生命安全；也为建筑施工装备的数字化、智能化升级提供了一种可落地、可推广的技术路径，具有重要的工程应用价值和现实意义。

\section{研究目的与意义}\label{ux7814ux7a76ux76eeux7684ux4e0eux610fux4e49}

\subsection{研究目的}\label{ux7814ux7a76ux76eeux7684}

针对塔吊吊装作业中吊钩端近场空间运动不确定、风险集中、环境干扰复杂等核心痛点，以及现有安全感知系统对吊钩端覆盖不足、单一传感器鲁棒性差、预警精度与实时性难以兼顾的问题，本研究旨在实现以下目标：

\begin{enumerate}
\def\labelenumi{(\arabic{enumi})}
\item
  设计并构建一套可工程化落地的吊钩端集成式感知系统，通过双激光雷达与工业相机的异构协同部署，以双雷达侧向对称布置实现近场盲区互补，以一体化盒体集成满足现场部署可靠性要求，突破传统传感器在覆盖范围、环境适应性上的局限，实现对吊钩周围 20 米内高危近场区域的全方位、高精度感知。
\item
  提出基于雷视融合的多源数据处理与信息融合方法，整合激光雷达的三维几何测距优势与视觉实例分割网络的语义识别能力，实现吊装构件的精准三维空间定位与双基准安全包络建模，以及动态危险目标的有效识别与短时轨迹跟踪。
\item
  建立融合三维静态净空评估与时域碰撞预测的动态风险分级预警模型，通过基于包络球与包络圆柱的双基准安全建模实现静态空间风险评估，并引入碰撞时间（TTC）与预测净空的协同判定机制实现动态风险升级，形成从感知、定位、判级到控制建议的完整预警闭环，提升预警的提前量与准确性。
\item
  验证系统在复杂施工工况下的稳定性、实时性与工程适用性，为塔吊吊装作业安全防护提供可推广、可复用的技术方案。
\end{enumerate}

\subsection{研究意义}\label{ux7814ux7a76ux610fux4e49}

（1）理论意义

现有塔吊安全感知研究多集中于塔身、塔臂等宏观层面，对吊钩端动态近场感知的研究较为匮乏。本研究提出“雷达为主、视觉为辅”的雷视融合架构，完善了多传感器在高危作业端的部署与标定方法，丰富了动态复杂场景下近场感知的理论体系。同时，研究建立的静动态目标分离、轨迹预测与分级预警模型，为多源数据融合在工程装备安全防护中的应用提供了新的思路，填补了吊钩端感知技术的研究空白。

（2）实践意义

本研究聚焦吊钩端这一最高风险区域，通过构建实时感知与预警系统，可减少人员误入、吊物碰撞等恶性事故发生概率，为施工现场安全管理提供技术保障。同时，系统采用轻量化硬件集成与高效算法设计，兼顾部署成本与运行稳定性，适配施工场地粉尘、震动、光照多变等恶劣环境，具备较强的工程落地性。

通过传感器集成、数据融合与智能算法的深度结合，本研究实现了吊装作业从“被动防护”向“主动预警”的转变，契合建筑行业安全升级与效率提升的核心需求。研究成果不仅可应用于塔吊设备，还可迁移至履带吊、汽车吊等其他起重机械，为工程机械行业的智能化转型提供典型示范。

\section{国内外研究现状}\label{ux56fdux5185ux5916ux7814ux7a76ux73b0ux72b6}

随着智能建造与施工现场数字化水平的不断提升，工程机械与大型吊装设备的安全感知与智能防护逐渐成为国内外学术界和工程界的研究热点。围绕施工设备运行安全，已有研究主要集中在作业环境信息获取以及碰撞风险识别与预警等方面，形成了从感知、理解到决策的多层次研究体系。然而，由于施工现场环境复杂、目标类型多样、作业过程动态性强，现有研究在感知对象、传感器布置位置以及安全防护粒度等方面仍存在一定局限，尤其在塔吊等大型起重设备的吊钩端动态作业区域，相关研究尚不充分。

\subsection{施工现场多源感知技术研究}\label{ux65bdux5de5ux73b0ux573aux73afux5883ux611fux77e5ux7814ux7a76}

施工现场环境感知是塔吊安全预警与避障控制的基础，其核心在于在复杂、动态且强不确定的工况下持续获取设备自身状态、作业空间几何以及人员与障碍物信息，并据此进行风险评估与预警。现有研究与工程应用中，塔吊相关感知传感器按功能定位大致可分为三类：编码器与IMU等本体传感器侧重于塔吊及吊钩端的运动状态与位姿获取，为感知预警提供统一参考基准；激光雷达、超声波等测距传感器提供环境几何距离与空间结构信息，服务于安全距离判断与入侵检测；视觉等语义传感器则聚焦目标类别识别，支撑对工人、车辆等危险目标的辨识与管理。三类传感器在测量维度、精度与环境鲁棒性方面各有优劣，决定其在塔吊避障中需根据工况进行针对性组合部署。然而，面向吊钩端近场场景，各类传感器在覆盖范围、环境适应性与融合策略上仍存在各自局限，以下就各传感器的技术特点与适用边界进行梳理。

（1）传感技术

施工现场常见的不同传感器在测量维度、精度与环境鲁棒性方面差异显著，工程实现通常采用多源组合以兼顾覆盖范围、成本与可维护性，如表\ref{tab:sensor_summary}所示。

\begin{table}[H]
  \centering
  \caption{施工现场常见传感器特性对比}
  \label{tab:sensor_summary}
  {\setlength{\tabcolsep}{3pt}
  \renewcommand{\arraystretch}{1.2}
  \renewcommand\tabularxcolumn[1]{m{#1}}
  \begin{tabularx}{\textwidth}{>{\centering\arraybackslash}m{2.0cm}>{\centering\arraybackslash}m{2.2cm}>{\centering\arraybackslash}X>{\centering\arraybackslash}X >{\centering\arraybackslash}m{2.8cm}}
    \Xhline{1pt}
    \textbf{传感器} & \textbf{信息维度} & \textbf{优势} & \textbf{局限} & \textbf{典型角色} \\
    \midrule
    编码器 & 回转角/幅度/高度 & 可靠、低成本、布置成熟 & 仅感知本体状态；外界障碍不可见 & 运动学/包络/禁行区输入 \\
    IMU & 角速度/加速度 & 动态姿态估计；运动畸变补偿支撑 & 漂移累积；需外部约束/融合 & 姿态与运动补偿；融合定位支撑 \\
    GNSS/RTK & 全局绝对坐标 & 开阔环境高精度基准 & 遮挡/多径；高动态精细支撑有限；不测障碍 & 全局参考与坐标对齐约束 \\
    UWB & 人员/标签位置 & 显式人员定位；电子围栏/入侵告警 & 依赖基站布设标定；金属/遮挡影响稳定性 & 人员风险提示 \\
    RFID & 构件身份 & 唯一标识；流程追溯/调度联动 & 无连续位置/距离；读写受干扰 & 身份/状态信息源 \\
    视觉（普通相机） & 语义 & 类别识别强；成本低、部署灵活 & 光照/遮挡敏感；精确测距能力弱 & 语义识别与可视化解释 \\
    视觉（深度/结构光） & 近距深度 & 近距测量；区域入侵检测 & 强光/远距/室外复杂工况受限 & 近距深度补充 \\
    毫米波雷达 & 测距 & 全天候；对运动目标敏感 & 角分辨率低；几何边界表达粗 & 冗余检测/存在性提示 \\
    超声波 & 近距测距 & 低成本；实现简单 & 量程短；易受介质/风噪干扰 & 近距防撞触发 \\
    LiDAR & 三维点云 & 稳定测距；不敏感光照；近场空间判定能力强 & 语义弱；反射/遮挡致缺失与离群点 & 几何测距与空间判定主传感器 \\
    \Xhline{1pt}
  \end{tabularx}
  }
\end{table}

编码器（包含角度、高度等状态传感器）常用于获取塔吊自身的状态，通过获取回转半径和回转角度等信息，能够逐步推算出小车位置、吊钩高度以及结构关键部位的运动边界。该类传感器可靠性高、成本低、工程布置成熟，能够为塔吊运动学模型、作业包络计算与禁行区生成提供直接输入，是典型的间接感知手段。然而其测量对象主要是设备自身状态，难以感知外部环境中临时出现的障碍物、人员靠近以及吊物周边的局部复杂结构，因此通常需要与外部环境传感器结合才能形成完整的避障能力。

惯性测量单元（IMU）能够提供角速度与加速度信息，在吊钩端强动态条件下用于运动状态估计、姿态变化监测以及对扫描类传感器（如LiDAR）进行运动畸变补偿。对于吊钩端避障场景，IMU的价值主要体现在两方面：其一，用于描述吊钩摆动、旋转等快速姿态变化，为点云稳定化与坐标系一致性提供支撑；其二，与其他传感器（如LiDAR、视觉、GNSS）融合形成更稳定的位姿估计。但IMU存在累计漂移，单独使用难以提供长期稳定的位置与姿态基准，通常需要与外部定位或地图约束结合。

GNSS（及RTK）常用于塔吊场景中的宏观定位与设备状态监测，例如塔吊基座或关键构件的绝对位置获取、施工区域级的空间约束建立等。RTK具备较高的平面定位精度，适用于开阔环境下的绝对坐标基准构建，也便于与BIM及施工组织规划进行空间对齐。然而在城市密集施工环境中，GNSS易受遮挡与多径效应影响，且对吊钩端等局部高动态部位的精细定位支撑有限；此外，GNSS本身无法直接感知障碍物几何边界，因此更适合作为全局参考或与IMU、LiDAR等共同构建多传感器定位框架。

UWB定位系统在施工安全管理领域应用较多，典型方式是在工人佩戴标签、现场布置基站，从而实现人员定位、电子围栏与区域入侵告警。UWB的优势在于可对人员提供显式定位信息，便于实现“人”的风险管理与统计，并且可与塔吊作业区域规则进行结合。但UWB系统部署依赖基站布设与标定，定位精度与稳定性受环境遮挡、金属结构、多径效应影响明显；同时UWB难以描述障碍物的几何形状与边界，通常更适合作为人员高危提示或辅助决策信息，而非替代基于几何测距的避障感知。

RFID技术在塔吊作业中更多用于构件、吊具或物料的身份识别与流程管理，例如对吊装构件进行唯一标识、实现装卸与流转记录、辅助吊装任务调度与追溯。在安全预警层面，RFID可用于识别“正在吊装的对象是什么”，与作业权限或工序规则联动，从而在管理层面增强安全性。但RFID通常难以提供稳定连续的空间位置与距离信息，其通信距离、读写可靠性也受金属遮挡与现场电磁环境影响，因此更适合作为“身份与状态信息源”融入系统，而非用于实时几何避障判定。

视觉作为一种信息丰富且成本低廉的信息获取方式，在语义理解方面具有明显优势，能够直观区分不同目标类别，且硬件成本较低、部署灵活。由于塔吊操作员的特殊工作位置导致不可避免存在视野盲区，早期研究者已尝试在塔吊上部署视觉系统\cite{KC3A3Q33}以显示吊装负载和下方工作区域，使驾驶室内的操作员能够直观看到吊钩端情况，减少视觉盲区。随着计算机技术发展，通过结合传统图像处理方法或深度学习算法，能够实现施工对象的目标检测，进而可以实现吊物空间姿态测量\cite{U6D69PTF}、生产效率统计\cite{MIA77CU2}、工人运动轨迹监测\cite{R625G22I}等任务。然而，视觉感知方法对光照条件和环境遮挡较为敏感，在夜间作业、强光或复杂遮挡条件下感知性能容易下降，同时难以直接获取精确的空间距离信息，在安全距离判断等应用中存在一定局限。与普通相机相比，深度相机及结构光相机能够直接提供局部深度信息，适合近距离空间测量与区域入侵检测，但在强光环境、远距离以及室外复杂工况下稳定性和测距范围会受到限制，其适用边界需要结合具体作业场景评估。

毫米波雷达具备全天候工作能力，对雨雾粉尘等恶劣工况鲁棒，且对运动目标具有较强敏感性，在夜间或能见度较低场景中具有优势。其不足在于角分辨率与空间细节表达能力通常不如LiDAR，难以精确刻画复杂障碍物的几何边界，因此在塔吊避障中更常用于冗余检测、近距离存在性判断或对运动目标的补充提示，而非单独承担精细空间建模任务。

超声波传感器成本低、实现简单，适合极近距离防撞与边界触发，如进行小范围防护或特定方向的距离阈值检测。但其测距范围有限，受传播介质、风噪与环境干扰影响较大，且难以覆盖吊钩端大范围三维空间，因此更多作为近距冗余或辅助触发传感器，与其他三维传感器配合使用。

激光雷达（LiDAR）能够输出三维点云并提供稳定的测距能力，是施工现场空间几何感知的关键传感器之一。对于塔吊避障而言，LiDAR可直接支持危险区域入侵判定、安全距离阈值触发、障碍物轮廓提取等任务，并且对光照变化不敏感，具备较高工程鲁棒性。但LiDAR对语义类别理解较弱，对场景物体的区分通常需要结合视觉或学习型算法；此外，金属反射、玻璃等材质以及遮挡引起的点云缺失、离群点等问题，需要配合滤波、聚类与一致性策略提高稳定性。因此在实际系统中，LiDAR常应用于空间测距、环境感知与建模等场景，与其他语义或状态传感器形成互补。

（2）多传感器融合感知

在复杂施工现场，仅依赖单一传感器往往难以同时满足安全预警所需的空间几何精度、语义理解能力与环境鲁棒性。因此，多传感器融合逐渐成为施工环境感知的重要技术路线。一般而言，多传感器融合是指在统一的时间与空间参考框架下，对来自不同传感器的信息进行对齐、关联与组合，以获得更完整、更可靠的环境状态估计或风险判定结果。其本质目标可概括为两类：其一，通过不同模态信息的互补性提升感知维度（例如视觉提供语义、测距传感器提供几何距离）；其二，通过信息的冗余性增强系统可靠性与容错能力，从而降低施工现场不确定因素带来的误报与漏报。

在塔吊避障与安全管理领域，多传感器融合研究大体可以分为两类路径：一类面向塔吊与吊钩端状态与位姿获取，以编码器、限位器、GNSS/RTK、IMU等为主要信息源，通过融合提升位姿连续性与全局参考稳定性，为作业包络计算、限制区生成与塔群协同防碰撞提供支撑。

许多研究者\cite{K9J9CXIH,NN865XZQ,8JDQ27IN,25XZJFDZ,5BYWKC7L,RBS3C7F5}明确了测量起重机操作的多个关键运动：吊臂回转角、提升绳伸长量和吊臂伸长量，并采用编码器、倾角传感器进行测量获取，以捕捉塔吊的运动状态。这种方式主要聚焦于布设于塔吊不同位置的传感器，根据每部分的参数通过建立塔吊的运动方程得到塔吊自身的状态，在原理上属于数据的拼接。仅依靠这些数据，在现实摆动的情况下无法准确估计吊钩和吊物的状态，这也是值得关注的高风险区域。对此，有通过线性激光装置测量塔吊小车到吊钩端的距离，再结合旋转编码器获取的旋转角度来得到吊物所在的位置\cite{RBS3C7F5}；也有在吊钩上侧金属壁安装IMU，并结合起重机的传感数据以捕捉吊钩的运动姿态（图~\ref{fig:transform-angular-measurements}）\cite{K9J9CXIH}。而Price等\cite{NN865XZQ}考虑到吊物冲击较大时可能损坏传感器，采用了视觉系统来定位吊物的偏移和旋转，但是缺乏吊物的深度信息，仍通过传感获取的缆绳长度作为吊物高度的间接估计。

\begin{figure}
  \centering
  \includegraphics[width=5.714in,height=1.666in]{image1.jpg}
  \caption{角度测量到绝对位置的映射示意\cite{K9J9CXIH}}
  \label{fig:transform-angular-measurements}
\end{figure}

由于施工现场是动态复杂环境，必须将现场环境数据集成到吊装辅助框架中才能起到实际的安全防护作用；同时，由于起重机特有的运动属性，其覆盖范围不能仅考虑二维，而是需要实现三维感知。三维激光扫描技术为上述需求提供了一种解决途径：依托地面激光扫描仪或无人机摄影测量技术\cite{G4AVEHSG}，可以获取工地范围内的3D点云数据，每个扫描点包含三维坐标信息及RGB颜色数据，可表征场景中物体的几何形状。扫描后的点云经过降采样、分割、聚类和方向估计等步骤，对场景中各种对象形成定向边界框，生成静态地图并集成到安全预警系统中，随周期性扫描进行更新\cite{NN865XZQ,K9J9CXIH}。但该方法仍存在分辨率与更新频率较低的问题。此外，扫描范围也受限于各塔吊的工作区域，大面积场景建模会产生大量冗余点云，为地图更新和算法判断带来处理性能压力。

另一类路径面向外部环境与危险目标风险识别，更强调对障碍物几何距离与目标语义信息的联合获取。典型做法包括以激光雷达/毫米波雷达提供可靠测距与空间入侵判定，以视觉方法提供危险目标（如工人、车辆等）的语义识别与可视化解释，并在决策层进行风险分级预警。在实际研究中，以避障为最终目标的研究通常使用多个传感器，在实现塔身姿态内部感知的基础上，完成外界状态的外部感知。

在施工现场环境感知中，视觉传感器因能提供目标语义信息与可视化解释，成为多传感器融合系统的核心感知单元，相关研究围绕技术优化与工程适配展开了大量探索。

Yang等\cite{8FY5FXQ3}采用 Mask R-CNN对人员与障碍物进行语义分割，通过像素 - 实际距离转换模型，联合输出人员语义状态与安全距离，计算误差控制在 3\% 以内；Golcarenarenji等\cite{MHXEXK6X}设计的 CraneNet深度学习模型，基于单目视觉实现 50 米范围内人员检测（语义识别）与远距离测量（几何距离），准确率达 92.59\%，且在嵌入式设备上实现 19 FPS的实时性，适用于起重机驾驶室等空间受限场景。这类方案的优势在于部署成本低、语义信息提取能力强，但测距精度受图像分辨率与特征提取质量影响，弱光环境下性能易下降。双目视觉方案则通过立体匹配直接获取障碍物距离，几何精度高于单目，同时结合 CNN实现人员与设备语义分类，但其测量范围受基线长度限制，近距离遮挡时匹配困难，设备体积也大于单目视觉，更适用于固定监测点的近距离作业场景。

视觉与其他传感器的融合方案进一步提升了联合感知的鲁棒性与精度。视觉与激光雷达融合方案中，视觉负责人员、设备、障碍物的语义识别，激光雷达提供高精度几何距离与三维点云，通过特征级融合将距离特征融入视觉检测网络，有效弥补视觉测距精度不足与激光雷达语义信息匮乏的缺陷，Price等\cite{T92DZ92V}基于该方案实现盲吊场景的障碍物距离与人员语义联合预警，适配高精度要求的复杂环境；视觉与RFID/IMU融合方案则利用 RFID/IMU的精准定位能力（几何距离）与视觉的行为语义提取优势，通过决策级融合实现人员动态风险预警，Zhang等\cite{G47UYPFK}采用 FairMOT算法跟踪人员与吊具，结合 Transformer模型预测轨迹，联合人员位置信息与移动语义，建立多等级碰撞风险预警规则，提升人员密集场景的风险识别可靠性。此外，聚焦人员风险识别的专项融合方案成为研究焦点，Shapira等\cite{KC3A3Q33}的视觉系统通过人员检测与距离判断，在缩短起重机作业周期的同时降低安全风险。

（3）雷视融合感知

在智能交通、自动驾驶、智慧工地等复杂环境感知场景中，精准、实时、鲁棒地获取目标信息是系统安全运行的核心前提。雷视融合通过整合激光雷达的三维几何测距与视觉的语义表达能力，可借助视觉语义优势弥补雷达短板，同时依托雷达空间感知能力修正视觉偏差，从而提升复杂场景下的感知鲁棒性与可解释性。

从融合策略与信息处理层级出发，现有研究将雷视融合划分为数据级融合、特征级融合、目标属性级融合与决策级融合四类\cite{PL2G5AIL}。

数据级融合强调在原始数据层面实现信息联合，对多源测距数据进行坐标统一与点云拼接、对图像与深度/点云进行时空对齐后再进行联合建模。该策略的关键在于统一不同模态数据的表达尺度与坐标基准，从而在最底层最大化保留信息量，并为后续处理提供更一致的输入表示。Zhou和Omar\cite{IPIVFMW7}采用像素级融合红外图像与RGB图像，采用自适应加权平均算法处理非饱和像素数据，针对饱和像素区域（眩光场景）实施主成分分析（PCA）算法，将两种输入的图像融合为单一输入，实现图像增强，以更鲁棒地检测实际场景中出现的人员和障碍。由于雷达和图像在水平与垂直分辨率上远低于光学图像\cite{PL2G5AIL}，因此，需要在时空间上分别进行对齐，而非图像间的归一化像素级融合。Kim等\cite{6VLQ22DN}综合使用 GPS、毫米波雷达、LiDAR与视觉构建环境表示，通过多帧 LiDAR观测累积生成网格地图，并以栅格“障碍观测计数”超过阈值触发风险预警，同时用雷达候选目标与视觉识别目标交叉验证并更新静态地图与安全区域。Lekic与 Babic\cite{6JL73W7F} 利用 GAN将雷达信息生成环境图像并与光学图像融合，Ouyang等\cite{9EZWD4IB}则使用条件 GAN在图像监督下由点云重建语义场景图像，旨在缓解点云与相机直接融合带来的计算负担，并在 KITTI数据集上验证实时检测有效性。

特征级融合通常先从各传感器数据中提取如几何特征、运动特征或深度网络特征的中间表示，再在特征空间完成联合建模，能够在信息利用与系统复杂度之间取得一定平衡，尤其在学习型感知框架中较为常见。传统机器学习时期往往需要先提取 HOG、GLCM等人工特征，再用 SVM与AdaBoost等分类器完成识别；而深度学习兴起后，特征融合更多被嵌入神经网络结构中，以端到端方式自动学习跨模态互补表示\cite{PL2G5AIL}。在雷视融合任务中，特征级融合常用于道路与可行驶区域检测等场景，例如Caltagirone等\cite{2GFRMDA6}基于全卷积网络实现 LiDAR与相机融合的道路检测。

更贴近工程避障预警需求的是目标属性级融合，各传感器先输出目标候选，再进行融合以降低单一传感器误报与漏报，该层级的抽象程度介于数据级融合与特征级融合之间，能够在保持计算可控的同时显著提升系统稳定性。一个与吊钩端避障可视化高度一致的范式是“LiDAR生成 ROI，再由图像中识别与解释 ROI”，如Wu等\cite{VWFQEPVD} 利用 LiDAR的距离与角度信息在图像中生成 ROI，并结合 3D点云形状信息进一步验证目标以降低误报、缓解遮挡带来的相机漏检。同时，该方法通过匹配 LiDAR与图像端的目标列表来提升检测速度并获得较高的行人检测精度。Han等\cite{B46NC5I5}从系统实现角度强调实时性与工程落地，为点云候选与视觉解释的在线管线提供了工程参考。类似地，Zhong等\cite{IPFD4E6N}利用毫米波雷达运动信息提供图像 ROI，并在 ROI上用 CNN识别目标，同时合并雷达与相机的目标列表以提升鲁棒性。

决策级融合则由各传感器独立完成检测、识别或风险判定，再通过规则、置信度或逻辑推理进行融合输出，如采用贝叶斯推理、D-S证据理论、模糊推理等。对工程避障而言，Wei等\cite{Z8QR9XPU}直接面向工业场景的实时避碰系统，体现了检测结果融合到避障决策的应用导向。De Silva等\cite{687GA9DL}则强调在移动机器人平台上实现鲁棒融合。

综上所述，雷视融合研究的核心共识在于：通过融合激光雷达的高精度三维几何与测距能力以及相机的纹理与语义表达能力，可以在复杂动态环境中显著提升目标感知的完整性、鲁棒性与可解释性，并降低单一传感器在遮挡、光照变化或传感噪声下造成的误检与漏检风险。现有工作虽然在融合深度与实现形式上差异较大，但普遍强调“可靠的时空对齐与外参标定”是雷视融合有效性的基础，并在此之上形成从端到端学习到工程化组合的多样路线。进一步而言，面向安全预警与避障等工程任务，雷视融合的研究趋势并非一味追求更深的网络耦合，而是更强调在实时性、系统复杂度与安全可靠性之间取得平衡：一方面利用几何信息实现稳定的距离与区域风险触发，另一方面用视觉语义增强对高风险目标的识别与可视化解释，从而构建“可落地、可维护、可扩展”的感知---预警闭环。

\subsection{吊装感知传感器布设方案研究}\label{ux540aux88c5ux611fux77e5ux786cux4ef6ux5e03ux7f6eux65b9ux6848ux7814ux7a76}

传感器布设方案（包括布设位置、传感器选型及集成方式）直接决定环境感知效果，其质量优劣也直接影响后续风险识别与避障决策的准确性与可靠性。为防范吊装作业对人员或建筑构件造成的损害，众多学者致力于通过智能感知技术与传感器部署来提升作业安全性，在塔吊吊臂、小车、驾驶室及施工现场周边等位置布设了摄像机\cite{FIQUTEC5}、惯性测量单元（IMU）\cite{NN865XZQ,FWEHHUWL}、射频识别（RFID）\cite{E97V5FCY}、超宽带（UWB）\cite{9LK3VTEQ}、超声波传感器\cite{TWWSPDKK}、高度传感器\cite{U3E4PYX6}、风速传感器\cite{2URRPIZ4}等多种传感器。根据传感器部署位置与感知方式的不同，这些研究可归纳为塔机本体间接测量、吊钩端直接传感与吊钩端视觉感知三类方案。

塔机本体布设方案将传感器安装于塔臂、小车、塔身等结构部位，通过多传感数据联立推算吊钩与吊物的空间位置，是当前研究中最为主流的部署方式。周飞虎\cite{USJ52UW5}在塔机回转机构、变幅机构与起升机构等关键部位布设角度、幅度与高度传感器，通过几何关系推算吊钩位置及构件相对关系，将潜在碰撞转化为构件间几何距离与阈值的比较以触发预警。Zhong等\cite{VNWJWSIA}构建了基于无线传感器网络（WSN）与物联网（IoT）的塔吊群安全管理系统，在塔臂、小车及起升机构等位置采集塔臂回转角度、小车位置、吊钩高度等信息，在统一坐标系下对塔吊结构空间关系进行计算以实现碰撞判断。针对塔吊间碰撞预警，Hwang等\cite{9LK3VTEQ}在每台塔吊吊臂前端和后端各布置1个UWB标签，通过AoA与TDoA相结合的方式获取标签空间位置，将安全距离映射为标签距离阈值。赵宇\cite{R5GR6N2A}沿吊臂方向均匀安装多个节点式超声波探测器阵列以形成线性保护区域，并融合高度与角度传感器推算吊装物体空间位置。Zhou等\cite{I2FL5SGE}将GNSS接收机固定于塔吊小车顶部，实时获取塔臂端部空间位置以监测多塔空间关系。张知田等\cite{XBDUPUTH}进一步引入视觉传感，在驾驶室下方与小车下方分别安装摄像头以捕捉工人分布与吊钩偏移，并结合塔吊主体传感器数据计算吊钩三维坐标，但该方法在夜间或扬尘工况下检测性能显著下降，在长吊臂（>30\,m）工况下像素级误差将被放大为分米级定位偏差\cite{XBDUPUTH}。总体而言，本体布设方案的传感器集中于塔机结构本体，工程部署难度低、成本可控，但都通过数据联立间接推算位置，而非直接感知吊钩端近场环境；在吊索摆动、遮挡与传感器视场限制下，对端部高风险区域的覆盖能力不足，且缺乏对周围障碍物及吊装作业直接风险的感知能力。

为克服本体布设方案对吊钩端感知不足的局限，部分研究将传感设备或可被可靠观测的靶标直接部署于吊钩块或吊物上，以更直接地获取端部真实运动状态并提升预警可靠性。Lee等\cite{J95BVZBC}在塔吊盲吊导航系统中于变幅臂端安装激光传感器、于吊钩块安装反射板作为测距靶标，利用激光束反射测得臂端到吊钩块的实际垂直距离，并以放绳长度编码器作为备份测量源。但该方案仅获取垂直距离这一单一维度信息，无法全面描述吊钩与吊物的三维位姿，且激光反射板依赖视线对准，在吊钩大幅摆动时极易丢失信号。Fang等\cite{K9J9CXIH}将无线IMU刚性安装于吊物或吊钩端载荷上，通过IMU测得吊物的欧拉旋转角，在已知吊绳长度并假设近似刚性的前提下，将姿态角转换为吊物的三维空间相对位置，实现对吊物摆动状态的在线监测。不过，这类端部直测方案要求在吊物上增加额外硬件，存在电池维护、现场粗放工况下的损伤风险及金属环境对无线链路干扰等工程约束\cite{K9J9CXIH}。Ku等\cite{J95BVZBC}进一步在机器人化塔吊系统中将3D LiDAR直接安装于吊钩块，获取近场点云以感知周围结构与潜在障碍，并在吊钩块上的工业计算机中实时提取路径方向上的障碍高度等关键信息。但该方案的点云处理目标仅局限于提取单一方向上的障碍高度，未充分利用LiDAR的全方位三维感知能力对吊物周边完整近场环境进行建模；同时未引入视觉传感器，缺乏语义识别能力，难以区分施工人员、车辆与普通障碍物等不同危险目标\cite{J95BVZBC}。

摄像机因其轻便、低成本且具备语义识别能力的优势，被研究者直接安装于吊钩端以实现对吊装下方区域的视觉感知，用于施工危险对象识别\cite{MHXEXK6X}、工人与负载碰撞预警\cite{R625G22I}等安全监测任务。然而，单一视觉方案在光照变化、扬尘遮挡等恶劣工况下性能显著下降，且缺乏深度信息，难以提供精确的三维空间定位。

综合上述三类方案可见，传感器部署呈现从塔机本体间接测量到吊钩端直接传感再到视觉感知的逐步靠近风险源的发展趋势，但现有方案在感知维度、环境鲁棒性与语义能力之间仍难以兼顾，亟需面向吊钩端场景构建融合激光雷达三维几何测量与视觉语义识别的多模态感知方案。

\subsection{碰撞风险监测与预警模型研究}\label{ux78b0ux649eux68c0ux6d4bux7b97ux6cd5ux4e0eux9884ux8b66ux6a21ux578bux7814ux7a76}

关于塔吊碰撞的检测与约束建模，既有研究在塔吊吊装规划与虚拟仿真场景中形成了较为系统的方法体系\cite{NEKKIG45}，总体呈现出从静态几何判别向动态风险评估、从二值碰撞判别向分级预警的演进趋势。根据建模对象与理论基础的差异，这些研究可归纳为几何空间驱动的碰撞检测与分级预警、时间域与预测驱动的风险建模、以及高效计算策略与动态避碰理论三类方法路径。

几何空间驱动的碰撞检测与分级预警是塔吊安全研究中最基本的建模方式，通过构建塔吊与障碍物之间的空间几何关系实现碰撞判别与风险触发。在虚拟仿真与路径规划场景中，研究人员多采用虚拟环境构建手段模拟实际吊装场景，以碰撞检测作为4D路径规划的约束条件\cite{KAVTM3MI}，通过为塔吊构建定向包围盒（Oriented bounding box，OBB）\cite{6SVWRM5Y}、采用离散碰撞检测（Discrete Collision Detection，DCD）\cite{HZ3G3PMB,3T2HAT9Q}或最小距离计算\cite{8FRKDLAT}等方式判断碰撞是否发生；Lai和Kang\cite{39QNQEMU}以简化外边界体近似现场机械与结构构件以降低计算负担。然而，这类方法依赖障碍物信息已知且完全可观测的假设，在实际场景中因传感器的观测范围、精度与延迟而难以直接迁移。针对施工过程的实时安全辅助，研究从“是否碰撞”的二值判别转向“碰撞风险强度”与“分级预警”建模，典型思路是构建安全包络或风险区域，采用最小距离与侵入判别触发告警\cite{KF34Y86G}。Wang等\cite{5JE9BPJQ}将风险区域与目标简化为AABB，通过进入/逼近判别实现安全预警；Yang等\cite{JSDZGAGY}使用Mask R-CNN检测危险源，根据图像像素与实际距离的转换关系计算安全距离。更进一步，面向吊装作业的实时主动安全辅助框架强调动态工作空间建模与风险提示机制\cite{K9J9CXIH}，Fang等\cite{FWEHHUWL}从人机协同角度讨论了安全风险呈现与操作者态势感知评估，为预警模型的工程落地提供了参考。

当障碍物与塔吊及吊物存在相对运动时，单纯依赖空间距离阈值进行判断存在固有局限，因此许多研究引入时域替代性风险指标，将风险刻画从“空间距离”推进到“时间余量”。典型如碰撞时间（Time-to-Collision，TTC）等指标，Nadimi等\cite{MAG8XRJ5}对TTC的适用性与局限进行了系统讨论；将TTC类思想与相对速度/加速度估计结合，可形成“距离—速度联合预警”框架，使预警更关注未来短时窗内的风险态势而非常态几何接触。进一步地，预测驱动的碰撞预警先对作业对象进行短时轨迹预测，再对未来时窗内的碰撞可能性进行判别或分级。Kim等\cite{XULW2BE3}指出预测潜在碰撞相较于仅预防邻近事故能更为有效地避免危险。张冬\cite{3FJ56MZB}将碰撞预警区分为基于距离的预警（针对静止建筑物）与结合速度模型的综合判断（针对塔吊间碰撞）。在轨迹预测方法层面，形成了从统计模型到深度学习模型的谱系，包括隐马尔可夫模型（HMM）\cite{TC4V9RAC}、长短期记忆人工神经网络（LSTM）\cite{HRDHH5IM}、深度神经网络（DNN）\cite{TF9IX4KK}等。在获得双方未来轨迹后，可通过碰撞锥\cite{UHXM7SNP,BHEQBQYN}或约束优化\cite{UYMXPTM6}等方法进行综合风险判断。

动态场景下的碰撞检测还面临实时性、精度与漏检风险之间的工程折中问题。为此，算法实现上通常采用“分层剔除与精细检测”策略：Dutta等\cite{HZ3G3PMB}在决策支持部分采用多层级OBB进行风险判别；Zhu等\cite{YMH5BX8M}强调融合空间划分与层次包围体并讨论了连续碰撞检测（CCD）在吊装扫掠体场景中的重要性；Lin等\cite{97BJ5NLW}提出基于点云的碰撞检测方法，反映出直接用现场环境状态驱动碰撞判别的趋势。在更一般的动态避碰理论中，碰撞风险还可通过相对速度可达集合进行刻画。碰撞锥与速度障碍（VO）范式可将塔吊端执行体与动态障碍的相对运动关系显式化\cite{BZUL7ZKB}，Van den Berg等\cite{25CQDW8W}将该思想扩展到多主体实时避碰。进一步地，安全关键建模将避碰问题表述为控制约束，控制障碍函数（CBF）成为近年来重要方向，Jian等\cite{QG5TNYGE}提出的D-CBF将障碍物预测与安全约束结合，为预警模型与后续控制策略的耦合提供了可迁移的建模范式。

上述研究为构建面向塔吊作业的实时碰撞预警系统提供了可复用的方法基础，也揭示了碰撞预警系统在实时性、准确性与鲁棒性间的固有矛盾，需在系统设计中加以权衡。

\subsection{文献综述小结}\label{ux6587ux732eux603bux7ed3}

塔吊高空作业环境与柔性吊索的动态特性，使上述各方向的方法在直接迁移至吊钩端场景时面临共性挑战。综合国内外研究现状，围绕塔吊吊装安全的环境感知方法、感知硬件布置、碰撞检测与预警模型三大核心方向，现有研究已形成一定的技术积累，但仍存在以下关键不足，为本文研究提供了明确的切入点：

（1）环境感知方法存在局限：单一传感器感知存在明显短板，如视觉受光照影响、雷达语义匮乏；多传感器融合成为主流趋势，但现有融合方案或过于依赖复杂深度学习模型导致实时性不足，或仅停留在数据拼接层面未能充分发挥互补优势。针对施工场景的雷视融合研究，尚未形成兼顾几何精度、语义识别与工程鲁棒性的成熟方案，尤其缺乏针对吊钩端动态特性的适配优化。

（2）感知硬件部署聚焦不足：现有传感器布置多集中于塔吊塔身、塔臂或施工现场固定位置，以宏观防碰撞为目标，对吊钩端这一最高风险区域的直接感知研究匮乏。虽有少数研究尝试在吊钩端部署传感设备，但存在抗干扰能力弱、感知范围有限等问题，难以满足近场全方位、高精度感知与跨模态数据融合需求。

（3）预警模型适配性不足：现有碰撞检测算法多源于虚拟仿真或路径规划场景，依赖静态几何模型，难以应对施工现场动态干扰多、障碍物复杂的实际情况；预警机制多基于单一空间距离阈值，未充分结合目标运动趋势与时间裕度（如TTC等时域指标），易出现误报或预警滞后问题。同时，对危险目标（如施工人员、车辆等）的语义关联不足，难以实现精准分级预警。

为直观对比现有方案与本文方法在各关键技术维度上的覆盖程度，表~\ref{tab:comparison_gap}从传感器安装位置、环境感知能力、目标识别与跟踪、碰撞预警策略及系统完整性等维度进行了系统梳理。

\begin{table}[htbp]
  \centering
  \newcommand{\cmark}{{\color{green!50!black}$\checkmark$}}
  \newcommand{\xmark}{{\color{red!60!black}$\times$}}
  \newcommand{\pmark}{{\color{orange!70!black}$\triangle$}}

  \small
  \renewcommand{\arraystretch}{1.6}
  \setlength{\tabcolsep}{6pt}

  \caption{现有方法与本文方法在吊钩端感知关键技术维度的对比}
  \label{tab:comparison_gap}

  \begin{tabular}{>{\centering\arraybackslash}m{3.6cm} *{6}{>{\centering\arraybackslash}m{1.3cm}}}
    \Xhline{1pt}
    &  % 留空，由第二行的 multirow 向上合并
    \multicolumn{3}{c}{\bfseries 大臂方案} &
    \multicolumn{3}{c}{\bfseries 吊钩端方案} \\
    \cmidrule(lr){2-4} \cmidrule(lr){5-7}
    \multirow{-2}{*}{\bfseries 对比维度} &
    \bfseries\makecell{编码器} &
    \bfseries\makecell{超声波/ \\ UWB} &
    \bfseries\makecell{激光\\雷达} &
    \bfseries IMU &
    \bfseries 视觉 &
    \bfseries\makecell{本文\\方法} \\
    \midrule
    安装位置靠近危险源   & \xmark & \xmark & \xmark & \cmark & \cmark & \cmark \\
    空间覆盖范围广        & \xmark & \xmark & \cmark & \xmark & \pmark & \cmark \\
    单帧感知无盲区        & \xmark & \xmark & \xmark & \xmark & \xmark & \cmark \\
    三维测距精度高        & \xmark & \pmark & \pmark & \xmark & \xmark & \cmark \\
    环境适应性强          & \cmark & \pmark & \cmark & \cmark & \xmark & \cmark \\
    具备语义识别能力      & \xmark & \xmark & \xmark & \xmark & \cmark & \cmark \\
    吊物状态估计精度高    & \pmark & \xmark & \pmark & \pmark & \xmark & \cmark \\
    动态目标跟踪能力强    & \xmark & \pmark & \pmark & \xmark & \pmark & \cmark \\
    碰撞预警时效性高      & \xmark & \xmark & \pmark & \xmark & \xmark & \cmark \\
    抗遮挡能力强          & \cmark & \pmark & \pmark & \cmark & \xmark & \cmark \\
    系统功能完整度高      & \xmark & \xmark & \xmark & \xmark & \xmark & \cmark \\
    \Xhline{1pt}
  \end{tabular}

  \parbox{\textwidth}{\vspace{0.5em}\small\raggedright 注：\cmark\;具备~~\xmark\;不具备~~\pmark\;部分具备}
\end{table}


\section{主要研究内容}\label{ux4e3bux8981ux7814ux7a76ux5185ux5bb9}

本文针对塔吊吊装作业中吊钩端近场高风险、环境遮挡多、光照变化大等问题，以激光雷达提供三维几何测距、视觉提供语义识别为基本思路，构建了一套在线碰撞风险判定与分级预警方法体系，具体研究内容包括以下四个部分：

（1）吊钩端感知系统构建与时空标定：设计双激光雷达与工业相机融合的吊钩端感知系统，完成一体化防护盒体及传感器集成；通过多传感器标定与时间同步建立统一时空基准，实现跨模态数据的空间对齐与时间一致性。

（2）基于融合点云的静态场景建图与动态目标感知：准确获取吊钩端位姿及近场静态场景，构建基于M-detector的动态目标检测与跟踪方法，识别施工人员等危险目标并估计其运动状态。

（3）基于雷视融合的吊物状态估计与三维定位：实现吊装构件的视觉识别，通过雷视融合将视觉语义映射至三维空间，实现吊装构件的空间定位；在此基础上构建静动态双基准安全包络模型。

（4）吊装作业风险分级预警与避障控制接口设计：建立基于静动态协同判据的风险分级预警机制，并设计与控制层解耦的标准化接口，形成从感知、定位、判级到控制建议的完整预警闭环。

在上述工作的基础上，本文各章安排如下：第1章阐述研究背景、意义与研究现状；第2章介绍系统总体架构与多传感器标定；第3章阐述基于融合点云的定位建图与动态目标感知方法；第4章阐述基于雷视融合的构件语义识别与安全包络建模；第5章建立分级预警策略与控制接口原型；第6章搭建实验平台并验证系统性能；第7章总结与展望。

\section{技术路线}\label{ux6280ux672fux8defux7ebf}

本文的技术路线如图~\ref{fig:tech_route}所示。在硬件部署与时空基准方面，通过双雷达与工业相机在防护盒体中的集成部署及多传感器联合标定与 PTPv2（Precision Time Protocol version 2）授时，建立跨模态时空基准，并基于“动态定位主链+静态感知支链”相分离的分层空间参考体系规范坐标发布。在感知定位环节，以双雷达外参实现点云统一融合与 LIO-SAM（LiDAR Inertial Odometry via Smoothing and Mapping）紧耦合定位，输出吊钩端位姿轨迹与静态点云地图，进而经 M-detector 动态判别与聚类跟踪获得近场动态目标的位置与速度。在视觉语义融合环节，以 YOLO26 实例分割网络实现构件像素级识别，通过点云-图像投影实现三维空间定位与双基准安全包络建模。在预警决策阶段，建立基于三维净空的静态分级与基于净空预测及 TTC 协同判定的动态升级机制，输出结构化预警等级与控制建议。

\begin{figure}[htbp]
  \centering
  \includegraphics[width=\textwidth]{技术路线图.pdf}
  \caption{本文技术路线图}
  \label{fig:tech_route}
\end{figure}

